\clearpage
\unitheader{}{Example question answer key}{}
Try the example questions before reading the answers below.

\paragraph{Postponing weight decay (p.~\pageref{ex:a2-wd-timing}).}
\textbf{Mean loss: $3.21$ for both two-update and four-update delays.}
Matching cumulative weight retention does not preserve the training trajectory,
but both delays give similar mean loss to decay at every update.

\paragraph{Widening a Kaiming model by copying (p.~\pageref{ex:a2-copied-width}).}
\textbf{(a) Option 3; (b) Option 2.}
For expansion factor $k$, use LR $.001/k$ for hidden matrices and readout,
and $.001$ for embeddings and normalization parameters.
With negligible Adam epsilon, scaling a gradient history by a positive constant
leaves its normalized update unchanged. Hidden and readout weights must stay
$1/k$ of the original weights, so their LRs scale by $1/k$; embeddings and
normalization gains remain unscaled copies. Finite precision and nonzero epsilon
allow small numerical differences.

\paragraph{Doubling batch size halfway through training (p.~\pageref{ex:a2-batch-switch}).}
\textbf{Policy A: change neither LR nor WD.}
Its mean loss is $2.93867$, compared with $2.94538$ for the next-best policy B.
Compensating for fewer weight-decay applications per token does not necessarily
improve final loss.
